\documentclass[12pt, spanish]{article}
\usepackage[utf8]{inputenc}
\usepackage[T1]{fontenc}
\usepackage{babel}
\usepackage{geometry}
\geometry{a4paper, margin=2.5cm}
\usepackage{hyperref}
\usepackage{listings}
\usepackage{xcolor}
\usepackage{enumitem}
\usepackage{graphicx}
\usepackage{float}

\definecolor{codegreen}{rgb}{0,0.6,0}
\definecolor{codegray}{rgb}{0.5,0.5,0.5}
\definecolor{codepurple}{rgb}{0.58,0,0.82}
\definecolor{backcolour}{rgb}{0.95,0.95,0.92}

\lstdefinestyle{mystyle}{
    backgroundcolor=\color{backcolour},
    commentstyle=\color{codegreen},
    keywordstyle=\color{magenta},
    numberstyle=\tiny\color{codegray},
    stringstyle=\color{codepurple},
    basicstyle=\ttfamily\footnotesize,
    breakatwhitespace=false,
    breaklines=true,
    captionpos=b,
    keepspaces=true,
    numbers=left,
    numbersep=5pt,
    showspaces=false,
    showstringspaces=false,
    showtabs=false,
    tabsize=2
}
\lstset{style=mystyle}

\title{\textbf{Informe de Modernización del Proyecto Crowd-TA-UVE}}
\author{Estiven Martínez}
\date{\today}

\begin{document}

\maketitle

\begin{abstract}
El presente informe documenta el proceso completo de modernización del proyecto \textit{``The Crowd Thinks Aloud''} (Crowd-TA-UVE), realizado entre el 18 y 19 de agosto de 2026. El proyecto fue migrado desde Django 3.0.8 y Python 3.8 a Django 5.1 y Python 3.13, eliminando dependencias obsoletas, adaptando la configuración para entornos Windows, creando la base de datos SQLite, y poniendo en funcionamiento el servidor de desarrollo. Se incluyen todas las etapas, desde la actualización del entorno hasta la fusión de ramas en Git y la puesta en marcha del flujo completo del proyecto.
\end{abstract}

\section{Introducción}
El proyecto \textit{Crowd-TA-UVE} fue desarrollado originalmente con Django 3.0.8 y Python 3.8, y dependía de varias librerías que hoy en día están obsoletas o son incompatibles con las versiones modernas de Python (3.13). El objetivo de esta modernización fue hacer que el proyecto pudiera ejecutarse en el entorno actual del desarrollador (Windows 11, Python 3.13) sin errores, manteniendo la funcionalidad original y preparando el código para futuras mejoras.

El proceso se dividió en 7 etapas, documentadas a continuación.

\section{Etapa 1: Actualización del Entorno y Dependencias}
\subsection{Creación del Entorno Virtual}
Se creó un entorno virtual aislado utilizando Python 3.13.0, asegurando que todas las instalaciones futuras no afecten al sistema global.

\begin{lstlisting}[language=bash, caption={Comandos para crear y activar el entorno virtual}]
cd C:\Users\ESTIVEN MARTINEZ\Desktop\Practica\crowd-ta-UVE\Source_Code
python -m venv venv
.\venv\Scripts\activate
\end{lstlisting}

\subsection{Actualización de Pip}
\begin{lstlisting}[language=bash, caption={Actualización de pip}]
python -m pip install --upgrade pip
\end{lstlisting}

\subsection{Instalación de Django 5.1}
\begin{lstlisting}[language=bash, caption={Instalación de Django 5.1}]
pip install django==5.1
\end{lstlisting}

\subsection{Instalación de Dependencias Modernas}
\begin{lstlisting}[language=bash, caption={Instalación de dependencias actualizadas}]
pip install django-crispy-forms>=2.1
pip install crispy-bootstrap5>=0.7
pip install django-cors-headers>=4.3
pip install django-extensions>=3.2
pip install django-countries>=7.6
pip install pillow>=10.0
\end{lstlisting}

\subsection{Reemplazo de \texttt{django-gsheets} por \texttt{gspread}}
La librería \texttt{django-gsheets} (y su dependencia \texttt{gsheets}) estaban abandonadas y no son compatibles con Django 5.1. Se decidió reemplazarlas por las librerías estándar para interactuar con Google Sheets:

\begin{lstlisting}[language=bash, caption={Instalación de nuevas librerías para Google Sheets}]
pip install gspread>=6.0
pip install google-api-python-client>=2.0
pip install google-auth-oauthlib>=1.0
\end{lstlisting}

\subsection{Generación de \texttt{requirements\_modern.txt}}
\begin{lstlisting}[language=bash, caption={Exportación de dependencias}]
pip freeze > requirements_modern.txt
\end{lstlisting}

\subsection{Gestión de Versiones con Git}
\begin{lstlisting}[language=bash, caption={Comandos Git ejecutados}]
git init
git add .
git checkout -b ramaEstiven
git commit -m "Estado inicial del proyecto antes de la modernización (ramaEstiven)"
git remote add origin https://github.com/estiradikal/crowd-ta-UVE
git push -u origin ramaEstiven
\end{lstlisting}

\section{Etapa 2: Actualización de \texttt{settings.py}}
Se modificó el archivo \texttt{settings.py} para hacerlo compatible con Django 5.1 y con el entorno Windows.

\subsection{Cambios realizados}
\begin{itemize}[label=--]
    \item \texttt{DEBUG = True} (para desarrollo local).
    \item \texttt{ALLOWED\_HOSTS = ['localhost', '127.0.0.1']}.
    \item Base de datos SQLite configurada.
    \item Rutas de archivos estáticos adaptadas a Windows.
    \item Se agregó \texttt{CorsMiddleware} al principio de \texttt{MIDDLEWARE}.
    \item Se agregó \texttt{'django\_language\_field'} a \texttt{INSTALLED\_APPS}.
    \item Se agregó \texttt{DEFAULT\_AUTO\_FIELD = 'django.db.models.BigAutoField'}.
    \item Se eliminó la opción obsoleta \texttt{USE\_L10N}.
\end{itemize}

\section{Etapa 3: Eliminación de Dependencias Obsoletas}
\subsection{Reemplazo de \texttt{django-language-field} por \texttt{CharField}}
El paquete \texttt{django-language-field} no se instalaba correctamente, por lo que se decidió reemplazar el campo \texttt{LanguageField} en los modelos por un \texttt{CharField} estándar.

\begin{itemize}[label=--]
    \item En \texttt{models.py}: \texttt{mother\_tongue = models.CharField(max\_length=32, null=False, default="unknown")}
    \item En \texttt{forms.py}: \texttt{mother\_tongue = forms.CharField(max\_length=32, required=True, label="Mother Tongue")}
    \item Se eliminó \texttt{'django\_language\_field'} de \texttt{INSTALLED\_APPS}.
\end{itemize}

\subsection{Eliminación de \texttt{gsheets.mixins}}
Las clases \texttt{TaskStatus} y \texttt{Approved\_Testers} ya no heredan de los mixins de \texttt{gsheets}. Se eliminó la dependencia completamente.

\begin{itemize}[label=--]
    \item Se comentó \texttt{from gsheets import mixins} en \texttt{models.py}.
    \item Se eliminó la herencia de \texttt{mixins.SheetPushableMixin} y \texttt{mixins.SheetSyncableMixin}.
    \item Se comentó la URL \texttt{path('', include('gsheets.urls'))} en \texttt{urls.py}.
    \item Se eliminó \texttt{'gsheets'} de \texttt{INSTALLED\_APPS}.
\end{itemize}

\section{Etapa 4: Migraciones y Base de Datos}
\subsection{Creación de la Base de Datos}
Se eliminó el archivo \texttt{db.sqlite3} (si existía) y se ejecutaron las migraciones para crear las tablas desde cero.

\begin{lstlisting}[language=bash, caption={Comandos para migrar la base de datos}]
del db.sqlite3
python manage.py makemigrations
python manage.py migrate usabilityTest
python manage.py migrate
\end{lstlisting}

\subsection{Poblado de Datos Iniciales}
Se crearon dos tareas en la tabla \texttt{TasksDescription} para poder probar el flujo completo del proyecto.

\begin{lstlisting}[language=python, caption={Código para insertar tareas de prueba}]
from usabilityTest.models import TasksDescription

TasksDescription.objects.create(
    task_id='1',
    task_description='Encuentra el curso "Global Warming Science" en la plataforma EDX.',
    task_valid_question='¿Cuál es el nombre de uno de los instructores?',
    valid_ans_options='["Prof. John Smith", "Dr. Jane Doe"]',
    answer='John Smith'
)

TasksDescription.objects.create(
    task_id='2',
    task_description='Encuentra cursos avanzados sobre "Climate change".',
    task_valid_question='¿Cuál es el nombre de uno de los cursos?',
    valid_ans_options='["Climate Change Science", "Energy Systems"]',
    answer='Climate Change Science'
)
\end{lstlisting}

\section{Etapa 5: Corrección de Errores en \texttt{views.py}}
\subsection{Descomentado de Variables Globales}
Se descomentaron las líneas que definían las variables globales \texttt{Total\_Tasks\_}, \texttt{Total\_Active}, \texttt{Task\_1\_Answer} y \texttt{Task\_2\_Answer} para corregir el error \texttt{NameError} en la vista \texttt{payment\_id}.

\begin{lstlisting}[language=python, caption={Variables descomentadas en views.py}]
Total_Tasks_ = TasksDescription.objects.all().count()
Total_Active = SubjectProfile.objects.all().count()
Task_1_Answer = TasksDescription.objects.get(task_id=1).answer
Task_1_Answer = Task_1_Answer.split("|")
Task_2_Answer = TasksDescription.objects.get(task_id=2).answer
Task_2_Answer = Task_2_Answer.split("|")
\end{lstlisting}

\subsection{Corrección de Importación de \texttt{redirect}}
Se agregó la importación de \texttt{redirect} en \texttt{urls.py} para corregir el error \texttt{NameError} en la redirección de la raíz.

\begin{lstlisting}[language=python, caption={Importación de redirect}]
from django.shortcuts import redirect
\end{lstlisting}

\section{Etapa 6: Configuración de Crispy Forms con Bootstrap 4}
El proyecto original usaba Bootstrap 4, pero se había instalado \texttt{crispy-bootstrap5}. Se corrigió instalando \texttt{crispy-bootstrap4} y ajustando \texttt{settings.py}.

\begin{lstlisting}[language=bash, caption={Cambio a Bootstrap 4}]
pip uninstall crispy-bootstrap5 -y
pip install crispy-bootstrap4
\end{lstlisting}

\begin{lstlisting}[language=python, caption={Configuración en settings.py}]
INSTALLED_APPS = [
    ...
    'crispy_forms',
    'crispy_bootstrap4',
    ...
]

CRISPY_ALLOWED_TEMPLATE_PACKS = ("bootstrap4",)
CRISPY_TEMPLATE_PACK = "bootstrap4"
\end{lstlisting}

\section{Etapa 7: Gestión de Ramas en Git y Fusión}
\subsection{Creación de la Rama \texttt{master}}
Se creó la rama \texttt{master} a partir de la rama remota \texttt{origin/master} (que ya existía en GitHub como rama por defecto).

\begin{lstlisting}[language=bash, caption={Creación de master local}]
git fetch origin
git checkout -b master origin/master
\end{lstlisting}

\subsection{Eliminación de la Rama \texttt{main}}
Se eliminó la rama \texttt{main} local y remota para mantener solo \texttt{master} como rama principal.

\begin{lstlisting}[language=bash, caption={Eliminación de main}]
git branch -D main
git push origin --delete main
\end{lstlisting}

\subsection{Fusión de \texttt{ramaEstiven} en \texttt{master}}
Se fusionaron todos los cambios modernizados de \texttt{ramaEstiven} en \texttt{master} usando la bandera \texttt{--allow-unrelated-histories} para permitir la fusión de historias diferentes.

\begin{lstlisting}[language=bash, caption={Fusión de ramaEstiven en master}]
git checkout master
git merge ramaEstiven --allow-unrelated-histories
git push origin master
\end{lstlisting}

\section{Estado Actual del Proyecto}
Tras completar todas las etapas, el proyecto se encuentra en el siguiente estado:

\begin{itemize}[label=--]
    \item \textbf{Entorno:} Python 3.13 con Django 5.1, entorno virtual \texttt{venv}.
    \item \textbf{Dependencias:} Todas las librerías actualizadas y compatibles.
    \item \textbf{Base de datos:} SQLite con migraciones aplicadas y datos iniciales (2 tareas).
    \item \textbf{Código:} Adaptado a Django 5.1, sin dependencias obsoletas (\texttt{gsheets}, \texttt{django-language-field}).
    \item \textbf{Git:} Rama \texttt{master} como rama principal, con todo el código modernizado.
    \item \textbf{Servidor:} Funcionando localmente en \texttt{http://127.0.0.1:8000/} con redirección a \texttt{/index/?workerId=test\&campId=test}.
    \item \textbf{Flujo:} Registro de usuarios, entrenamiento, tareas y pago funcionando correctamente.
\end{itemize}

\section{Próximos Pasos}
\begin{itemize}[label=--]
    \item Probar el flujo completo en diferentes navegadores.
    \item Ajustar el campo \texttt{Mother Tongue} para que tenga opciones predefinidas (si se desea).
    \item Implementar la sincronización con Google Sheets usando \texttt{gspread} (actualmente comentada).
    \item Configurar variables de entorno para secretos y claves de API.
    \item Preparar el proyecto para producción (cambiar \texttt{DEBUG = False}, configurar servidor WSGI/ASGI).
\end{itemize}

\section{Conclusiones}
La modernización del proyecto \textit{Crowd-TA-UVE} se ha completado con éxito. El proyecto ahora es ejecutable en el entorno actual del desarrollador con Django 5.1 y Python 3.13, sin errores de sintaxis o dependencias obsoletas. La creación de una rama \texttt{master} estable permite que otros colaboradores puedan clonar y trabajar sobre el código sin problemas. El flujo completo del proyecto (registro, tareas, pago) ha sido probado y funciona correctamente en el entorno de desarrollo local.

\end{document}